\documentclass[11pt]{article}
\usepackage[margin=1in]{geometry}
\usepackage{amsmath,amssymb,amsfonts,mathtools}
\usepackage{array,booktabs,enumitem,graphicx,xcolor,hyperref}
\usepackage[T1]{fontenc}
\usepackage[utf8]{inputenc}
\title{How to Put a Heavier Higgs on the Lattice}
\author{U.M. Heller, H. Neuberger and P. Vranas}
\date{}
\begin{document}
\maketitle
\begin{abstract}

Lattice work, exploring the Higgs mass triviality bound, seems to indicate that a strongly interacting scalar sector in the minimal standard model cannot exist while low energy QCD phenomenology seems to indicate that it could. We attack this puzzle using the 1/N expansion and discover a simple criterion for selecting a lattice action that is more likely to produce a heavy Higgs particle. Our large \$N\$ calculation suggests that the Higgs mass bound might be around \$850 GeV\$, which is about 30\% higher than previously obtained.

\end{abstract}

\section{Introduction}

Lattice work, exploring the Higgs mass triviality bound, seems to indicate that a strongly interacting scalar sector in the minimal standard model cannot exist while low energy QCD phenomenology seems to indicate that it could. We attack this puzzle using the math expression expansion and discover a simple criterion for selecting a lattice action that is more likely to produce a heavy Higgs particle. Our large math expression calculation suggests that the Higgs mass bound might be around math expression , which is about 30\textbackslash{}\% higher than previously obtained.

\section{Method}

ACKNOWLEDGEMENTS. This research was supported in part by the Florida State University Supercomputer Computations Research Institute which is partially funded by the U.S. Department of Energy through DOE grant \textbackslash{}\# DE-FG05-85ER250000 (UMH and PV) and under DOE grant \textbackslash{}\# DE-FG05-90ER40559 (HN). UMH also acknowledges partial support by the NSF under grant \textbackslash{}\# INT-8922411 and he would like to thank F. Karsch and the other members of the Fakult\textbackslash{}"at f\textbackslash{}"ur Physik at the University of Bielefeld for the kind hospitality while part of this work was done. HN would like to thank F. David and M. Einhorn for useful discussions.

\begin{equation}
\label{eq:compact}
L(\theta) = \sum_{i=1}^{n} \ell(x_i, \theta)
\end{equation}
Equation~\ref{eq:compact} is included to keep the sample paper-like while remaining small.

\section{Conclusion}

Lattice work, exploring the Higgs mass triviality bound, seems to indicate that a strongly interacting scalar sector in the minimal standard model cannot exist while low energy QCD phenomenology seems to indicate that it could. We attack this puzzle using the math expression expansion and discover a simple criterion for selecting a lattice action that is more likely to produce a heavy Higgs particle. Our large math expression calculation suggests that the Higgs mass bound might be around math expression , which is about 30\textbackslash{}\% higher than previously obtained.
\end{document}
